%%%%%%%%%%%%%%%%%%%%%%%%%%%%%%%%%%%%%%%%%%%%%%%%%%%%%%%%%%%%%%%%%%%%%%
%
% Example of \TeX Format for AQIS Pre-Proceedings
%
%%%%%%%%%%%%%%%%%%%%%%%%%%%%%%%%%%%%%%%%%%%%%%%%%%%%%%%%%%%%%%%%%%%%%%%

\documentclass[a4paper]{article}
\usepackage{aqis}

\begin{document}

\title{
   Example of \TeX\ Format for AQIS Pre-Proceedings
}

\author{
  David Clarino
  \affiliation{1}
  \email{username1@domainname1}
  \and
 Atsushi Matsuo
   \affiliation{2}
  \email{username2@domainname2}
  \and
 Shigeru Yamashita
   \affiliation{1}
  \email{username3@domainname3}
  }
  
\address{1}{
  Ritsumeikan University, Graduate School of Information Science and Engineering
}
\address{2}{
  IBM Research, Tokyo
}


\abstract{
  Lookup Table (LUT) based synthesis methods have recently been proposed as a way to synthesize quantum Boolean circuits
  in a qubit-constrained environment. Other recent research has demonstrated the possibility of using {\it relative phase}
  quantum circuits when compute/uncompute logic is used in tandem, reducing T-count in quantum Boolean circuits in
  the fault-tolerant quantum computing paradigm. Because LUT-based synthesis methods use compute/uncompute pairs on ancilla
  qubits, this suggests that implementing the arbitrary Boolean logic that make up the individual Boolean logic network nodes
  in a relative phase manner could reduce the T-count. To generate such arbitrary Boolean functions, we utilize Shannon's
  decomposition, Davio expansions, as well as alternating balanced and unbalanced relative  phase circuits. Experimental results
  demonstrate that our method can reduce the T-count to an average of 24\% of the existing method. 
}

\keywords{relative-phase Toffoli gates (RTOF), circuit optimization, LUT, quantum circuit synthesis}


\maketitle

%########### body

\section{Introduction}
{\it Quantum Boolean circuits}~\cite{bib-ddmf} implement Boolean operations as quantum circuits and
are common components of many quantum algorithms such
as oracles in Grover's algorithm and cost functions. To realize such quantum Boolean circuits,
a Toffoli gate, which implements the Boolean AND, is an essential logic primitive in a universal gate
set. These Toffoli gates have to be, in turn, composed of physically realizable gates, such as NOT,
CNOT, and T gates (e.g., the Clifford + T~\cite{bib-barenco-elementary} basis gate set), as depicted
in Fig.~\ref{fig-toffoli}. However, in the fault-tolerant paradigm, T-gates incur much higher cost
than the NOT and CNOT gates. It is therefore useful to consider ways to reduce this {\it T-count},
which is the number of both the T-gates and their inverse in the circuit.

One of the methods that has been proposed to reduce the T-count is to implement Boolean functions only
up to a {\it relative phase}~\cite{bib-barenco-elementary}, which is an input dependent phase on the
output. Some recent work~\cite{bib-rtof-maslov,bib-selinger-rtof,bib-amy-phase-state} has proposed
relative phase versions of common Boolean functions, which achieve T-count savings by removing the
T gates that only affect the phase.

Recent attempts to synthesize quantum circuits using LUT-based methods have shown some promising results
~\cite{bib-lhrs-soeken}. As intermediate nodes {\it cascade} their outputs into output nodes, logic produced
by such a method is exactly the type that can be improved by synthesizing relative phase versions of its
component logic. In addition, the fact that relative phase versions can be implemented without any additional
ancilla means that it works well to synthesize Boolean functions in the qubit constrained environment that
such methods are intended for. 

In this work we attempt to 
\paragraph{Basic Format:}

\begin{itemize}
  \item
      When you create PDF file from your manuscript source
      (especially from MS-Word), please select fonts carefully for
      some fonts cannot be handled well by PDF creating system.
  \item
    The manuscript should
    not exceed three pages including the title page,
    using A4-sized papers.
  \item
    No page of the manuscript should contain page numbers/page headings.
\end{itemize}

\paragraph{Layout:}

\begin{itemize}
  \item
    The manuscript should
    have top-, foot-, and side-margins of 2cm, 2.5cm, and 1.5cm,
    respectively (the title page should have top margin of 3cm).
  \item
    The manuscript should begin with a title,
    followed by names, affiliations and postal addresses of authors,
    followed by a 100-word abstract and keywords.
    The main body should use two-column style.
    E-mail addresses of authors should be placed in footnotes.
\end{itemize}

\section{Theorems and Proofs}

Our style file prepares the following environments.

\begin{center}
\begin{tabular}{ll}
\hline
Environment name & Heading\\
\hline
{\tt definition} & {\bf Definition}\\
{\tt theorem} & {\bf Theorem}\\
{\tt lemma} & {\bf Lemma}\\
{\tt corollary} & {\bf Corollary}\\
{\tt proposition} & {\bf Proposition}\\
{\tt example} & {\bf Example}\\
{\tt remark} & {\bf Remark}\\
{\tt fact} & {\bf Fact}\\
{\tt conjecture} & {\bf Conjecture}\\
\hline
\end{tabular}
\end{center}

\section{Figures and Tables}

For figures and tables, the captions should be placed
below figures and above tables.

\section{Bibliography Styles}

The following bibliography style used in this document is preferable.
The authors may use common abbreviations for journal and conference names.

\begin{thebibliography}{9}

\bibitem{Grover96}
L.~K.~Grover.
\newblock A fast quantum mechanical algorithm for database search.
\newblock In {\em Proc.\ of the 28th ACM STOC}, pages 212--219, 1996.

\bibitem{Gruska99}
J.~Gruska.
\newblock {\em Quantum Computing}.
\newblock McGraw-Hill, 1999.

\bibitem{Holevo96}
A.~S.~Holevo.
\newblock The capacity of quantum channel with general signal states.
\newblock quant-ph/9611023, 1996.

\bibitem{Sho97SIComp}
P.~W.~Shor.
\newblock Polynomial-time algorithms for prime factorization
and discrete logarithms on a quantum computer.
\newblock {\em SIAM J. on Comp.}, 26(5):1484--1509, 1997.

\end{thebibliography}

\end{document}
% end of file
